\begin{cnabstract}{实体组件系统；类吸血鬼幸存者；系统设计与实现；LayaAir；数据驱动}
随着 HTML5 与 WebGL 技术的普及，浏览器端实时动作游戏在架构解耦与同屏性能方面面临新的挑战。类吸血鬼幸存者玩法融合 Roguelite 单局构筑与元游戏流程，要求客户端在数据驱动、模块划分与可验证性上具备清晰边界。本文以 Survivor And Fight 毕业设计为研究对象，基于 LayaAir 3.x 与 TypeScript，对 2D 俯视角生存战斗游戏开展从需求分析、总体设计、系统实现到测试验证的完整研究。

论文主体共分八章：绪论与相关技术奠定选题依据与技术路线；需求分析明确功能、非功能需求及验收映射；总体设计提出五层架构，并从 ECS Gameplay、MVC 界面、对象池与 Web Worker 三方面给出方案；实现与性能章节分别论述核心模块、主循环及优化实践；测试章节通过单元脚本、功能用例与压力实验进行验证；结论章节归纳成果与不足。原型可稳定支撑菜单、跑图、战斗与技能装配等闭环流程，对象池与动态图集可改善高同屏负载下的运行表现。
\end{cnabstract}
